
\section{Hyperbolic parallelograms}
\label{sec:hyp-par}

\DPTtodo[inline]{Just putting this here for now, I don't know where it
  actually belongs in the paper}
  \DMGtodo[inline]{I think leaving it here is fine}

\begin{definition}
  A \emph{hyperbolic parallelogram} is a hyperbolic quadrilateral
  where opposite sides have equal length, as shown in
  Figure~\ref{fig:hyp-par}. Equivalently, it is a quadrilateral where
  the diagonals bisect each other.
\end{definition}

\begin{figure}
  \[
    \includegraphics{hyp-quad-0}
  \]
  \caption{A hyperbolic parallelogram}
  \label{fig:hyp-par}
\end{figure}

\DPTtodo{Figure should have more labels. But almost certainly the
  naming of the variables in this lemma should change, since $a$
  should be the name for the translation from $X$ to $Y$ along
  the line $XY$.}

% The following geometric lemma is presumably standard. It is the
% hyperbolic version of the parallelogram law that, in a parallelogram
% with sides $c,d$ and diagonals $2a, 2b$, we have $2(a^2+b^2) =
% c^2+d^2$. It can also be thought of as a formula giving the median of
% a hyperbolic triangle in terms of its side lengths.
\begin{lemma}[Hyperbolic parallelogram law]
  Given a hyperbolic parallelogram $XYZW$ with center $O$, let
  $c = \abs{X-Y} = \abs{Z - W}$ and $d = \abs{Y - Z} = \abs{X - W}$ be
  the lengths of the sides, and let $a = \abs{X - O} = \abs{Z
    - O}$ and $b = \abs{Y - O} = \abs{W - O}$ be the half-lengths of
  the diagonals. Then
  \[
    2 \cosh a \cosh b = \cosh c + \cosh d.
  \]
  \label{lem:hyp_parallelogram}
\end{lemma}

\begin{proof}
  Let $\gamma$ be the measure of $\angle XOY$ and $\delta=\pi-\gamma$ be the
  measure of $\angle YOZ$. The hyperbolic law of cosines gives
  \begin{align*}
    \cosh c &= \cosh a \cosh b - \sinh a \sinh b \cos \gamma\\
    \cosh d &= \cosh a \cosh b - \sinh a \sinh b \cos \delta.
  \end{align*}
  Adding these gives the desired result.
\end{proof}

%%% Local Variables:
%%% mode: latex
%%% TeX-master: "Intersections"
%%% End:
